\documentclass[11pt]{article}

\usepackage{amsmath,amssymb,amsthm}
\usepackage[margin=1in]{geometry}
\usepackage{enumitem}
\usepackage{mathtools}
\usepackage{booktabs}
\usepackage{graphicx}

\title{Preepi, Closed Convex Functions, and Linear Logic:\\
\large A Research Program}
\author{Thu-Le Tran\\
Can Tho University\\
\texttt{ttle@ctu.edu.vn}}
\date{}

\newtheorem{claim}{Claim}
\newtheorem{question}{Open Question}

\newcommand{\R}{\mathbb R}
\newcommand{\Rp}{\mathbb R_{+\infty}}
\newcommand{\epi}{\operatorname{epi}}
\newcommand{\dom}{\operatorname{dom}}
\newcommand{\cl}{\operatorname{cl}}
\newcommand{\conv}{\operatorname{conv}}
\newcommand{\fibsum}{+_{\mathrm{fib}}}
\newcommand{\Msum}{+_{\mathrm M}}
\newcommand{\Fact}{\mathsf{Fact}}
\newcommand{\parr}{\mathbin{\text{\rotatebox[origin=c]{180}{\&}}}}
\newcommand{\ri}{\operatorname{ri}}
\providecommand{\llbracket}{[\![}
\providecommand{\rrbracket}{]\!]}

\begin{document}
\maketitle

\begin{abstract}
We propose \emph{preepi} — arbitrary subsets of $X\times\R_{+\infty}$,
$\R_{+\infty}:=\R\cup\{+\infty\}$ — as a common geometric object linking
three worlds: closed convex functions on $X=\R^n$, a lattice-and-monoid
algebra closed under a Fenchel-type polarity, and the connectives of
multiplicative-additive linear logic (MALL), the fragment of linear logic
without exponentials. \textbf{Every statement in this document, including
every definition, is a claim: a proposal for how these objects should be
modeled and related, not an established result.} Definitions are
themselves modeling choices and are labeled as such; no distinction in
confidence is intended between a claim that happens to look like a
definition and a claim that happens to look like a theorem. We give the
proposed correspondence table, the proposed definitions, and a roadmap of
three milestones, deliberately separated by how much is claimed at each:
intrinsic properties of preepi, the connection to convex optimization, and
the connection to logic — the last of which is stated only as open
questions.
\end{abstract}

\tableofcontents

\section{Introduction}

Multiplicative-additive linear logic (MALL) is the fragment of Girard's
linear logic obtained by dropping the exponential modalities $!,?$ (and any
first- or second-order quantifiers): it retains exactly the four binary
connectives $\otimes,\parr,\&,\oplus_{\mathrm{MALL}}$, their units, and
linear negation $(\cdot)^\bot$. Full linear logic recovers resource
duplication and discarding via $!,?$ on top of this fragment; the present
program restricts itself to MALL throughout, since preepi are one-shot
geometric objects with no evident candidate structure for $!,?$ (see
Milestone III, Section~\ref{sec:milestones}).

The classical bridge between convex analysis and phase semantics for linear
logic runs, informally, as follows: a closed proper convex function
$f:X\to\Rp$ has an epigraph that behaves like a fact of a Girard quantale
under Fenchel conjugation, and the resulting algebra of facts is a
candidate poset-level model of MALL. This program proposes to make that
bridge precise by working entirely at the level of \emph{preepi} — raw
subsets of the extended space $X\times\Rp$ — treating the closed convex
function and the epigraph as two faces of a single object, and treating the
Fenchel polar as the negation of a phase-semantic pole.

We propose the following correspondence as the organizing table of the
program. Columns left to right: the MALL connective, the proposed preepi
operation, and the function-level operation it is meant to mirror.

\begin{center}
\begin{tabular}{lll}
\toprule
MALL & Preepi operation & Function-level mirror \\
\midrule
pole / pairing & $\beta(x,y):=\langle x,y\rangle$ & --- \\
$(\cdot)^\bot$ & $E\mapsto E^*$ (Fenchel polar) & $f\mapsto f^*$ \\
fact ($A=A^{\bot\bot}$) & $E=E^{**}$, i.e.\ $E\in\bar X$ & $f=f^{**}$ \\
$\vdash$ / provability order & $E\subseteq F$ & $f\ge g$ (\emph{note the reversal below}) \\
$\&$ (with) & $E\wedge F:=E\cap F$ & $\max(f,g)$ \\
$\oplus_{\mathrm{MALL}}$ (plus) & $E\vee F:=(E\cup F)^{**}$ & $(\min(f,g))^{**}$ \\
$\otimes$ (tensor) & $E\oplus F:=(E\fibsum F)^{**}$ & $f+g$ \\
$\parr$ (par) & $E\boxdot F:=(E^*\fibsum F^*)^*$ & $\cl(f\,\square\,g)$ \\
unit of $\otimes$ & $1_\oplus:=X\times[0,+\infty]$ & $\cl(0)$ \\
unit of $\parr$ & $1_\boxdot:=(\{0\}\times[0,+\infty])\cup((X\setminus\{0\})\times\{+\infty\})$ & $\cl(\delta_{\{0\}})$ \\
unit of $\&$ & $1_\wedge:=X\times\Rp$ & $f\equiv-\infty$ (excluded, see \S2) \\
unit of $\oplus_{\mathrm{MALL}}$ & $1_\vee:=X\times\{+\infty\}$ & $f\equiv+\infty$ \\
\bottomrule
\end{tabular}
\end{center}

\emph{On the $\otimes/\parr$ assignment.} The table assigns
$\otimes\leftrightarrow\oplus$ (fiber sum, mirroring $f+g$) and
$\parr\leftrightarrow\boxdot$ (mirroring closed infimal convolution), on the
resource-combination intuition that $\otimes$ should mirror the more
primitive, unconditionally-defined operation. This assignment is
\textbf{not} forced by the definitions alone: Girard's phase semantics
defines $\otimes$ directly from the single ambient monoid of the ground
space, and on $\tilde X$ the only operation satisfying that role
unconditionally is $\Msum$ (Minkowski sum), which would instead favor the
opposite assignment $\otimes\leftrightarrow\boxdot$, $\parr\leftrightarrow\oplus$.
Whether the assignment in the table survives, is reversed, or requires a
genuinely different ambient structure (e.g.\ a quantaloid rather than a
single quantale) is left open; see Milestone III,
Question~\ref{q:swap}. Any later revision of this table must re-check the
order-direction row above as well, since swapping $\otimes\leftrightarrow\parr$
also swaps which unit pairs with which connective.

\emph{On the order row.} $E\subseteq F$ is proposed as the preepi mirror of
provability, but at the function level the corresponding relation is
\textbf{reversed}: $f\ge g$ pointwise corresponds to $\epi(g)\subseteq\epi(f)$,
not $\epi(f)\subseteq\epi(g)$ — a larger function has a \emph{smaller}
epigraph. Every claim in Section~\ref{sec:components} that mixes the
preepi order with the function order must be checked against this reversal
individually; it is not safe to assume a naive order-preserving translation
between the two columns.

\section{The components}\label{sec:components}

Every item below is a \textbf{Claim}: a proposal for how the corresponding
object should be modeled. Motivation from the function-level picture is
given alongside each one.

\subsection{Preepi}

\begin{claim}[Ground space]
Fix
\[
X:=\R^n,\qquad \Rp:=\R\cup\{+\infty\},\qquad \tilde X:=X\times\Rp.
\]
A \textbf{preepi} is an arbitrary subset $E\subseteq\tilde X$; no separate
symbol is introduced for the family of preepis.
\end{claim}

\emph{Motivation.} A point $(x,r)\in\tilde X$ records a location together
with a height. The value $+\infty$ is retained, and $-\infty$ excluded, so
that a preepi can model a \emph{total} function $f:X\to\Rp$ — the standard
codomain of convex analysis — directly, without needing to encode
$f(x)=+\infty$ indirectly through an empty fiber.

\subsection{Four closure operators}

\begin{claim}[Closures]\label{cl:closures}
For $E\subseteq\tilde X$ and $x\in X$, write $E_x:=\{r\in\Rp:(x,r)\in E\}$.
Define
\begin{align*}
U(E) &:= \{(x,r'):\exists r\le r',\ (x,r)\in E\} &&\text{(upward closure)}\\
F(E) &:= \{(x,r): r\in\cl(E_x)\} &&\text{(fiber closure)}\\
C(E) &:= \conv(E) &&\text{(convex closure)}\\
T(E) &:= \cl(E) &&\text{(topological closure)}
\end{align*}
where $\cl$ and $\conv$ are taken with respect to the order topology on
$\Rp$ (in which $r_k\to+\infty$ converges) and the convex-combination rule
that declares a segment to run through $+\infty$ as soon as one endpoint
does.
\end{claim}

\emph{Motivation.} Each operator repairs exactly one property separating a
general preepi from the extended epigraph of a proper, convex,
lower-semicontinuous function: upwardness in each fiber, fiber closedness,
global convexity, global closedness.

\subsection{Regepi}

\begin{claim}[Regepi]\label{cl:regepi}
A preepi $E$ is a \textbf{regepi} if
\[
U(E)=F(E)=C(E)=T(E)=E.
\]
Write $\bar X:=\{E\subseteq\tilde X:\ E\text{ is a regepi}\}$.
\end{claim}

\emph{Motivation.} A regepi is proposed as exactly the class of preepis
arising as extended epigraphs of proper, convex, lower-semicontinuous
functions $X\to\Rp$.

\subsection{Fenchel polar and the order}

\begin{claim}[Polar and order]\label{cl:polar}
For $E\subseteq\tilde X$,
\[
E^*:=\{(y,s)\in\tilde X:\ \langle x,y\rangle\le r+s\ \ \forall(x,r)\in E\},
\]
with a fixed Euclidean inner product on $X$ and $r+(+\infty):=+\infty$.
The order on preepis is inclusion, $E\subseteq F$.
\end{claim}

\emph{Motivation.} This is the geometric form of the Fenchel conjugate
$f^*(y):=\sup_x\{\langle x,y\rangle-f(x)\}$: proposed so that
$E=\epi(f)\ \Rightarrow\ E^*=\epi(f^*)$. Because $s=+\infty$ trivially
satisfies the defining inequality, $E^*\ne\varnothing$ for every preepi
$E$ — a structural feature not available on $X\times\R$. As noted in
Section~\ref{sec:components} above (order row), $E\subseteq F$ mirrors
$f\ge g$, not $f\le g$, when $E=\epi(f)$, $F=\epi(g)$; this reversal must
be tracked, not assumed away, in any later claim comparing the two orders.

\subsection{Four operators}\label{sec:operators}

\begin{claim}[Primitive operations]\label{cl:primitive}
For $E,F\subseteq\tilde X$,
\[
E\fibsum F:=\{(x,r+s):(x,r)\in E,(x,s)\in F\},\qquad
E\Msum F:=\{(x_1+x_2,r_1+r_2):(x_1,r_1)\in E,(x_2,r_2)\in F\}.
\]
\end{claim}

\emph{Motivation.} $\fibsum$ mirrors pointwise addition of functions,
$\epi(f)\fibsum\epi(g)\leftrightarrow\epi(f+g)$, with no closure needed at
the function level. $\Msum$ mirrors the numerator of infimal convolution,
$\epi(f)\Msum\epi(g)\leftrightarrow\epi(f\,\square\,g)$ before closure.

\begin{claim}[Induced operations]\label{cl:induced}
For $E,F\in\bar X$,
\[
E\wedge F:=E\cap F,\qquad
E\vee F:=(E\cup F)^{**},\qquad
E\oplus F:=(E\fibsum F)^{**},\qquad
E\boxdot F:=(E^*\fibsum F^*)^*.
\]
\end{claim}

\emph{Motivation.} $\wedge,\vee$ mirror $\max(f,g)$ and the bipolar closure
of $\min(f,g)$. $\oplus$ mirrors closed pointwise addition, $f+g$.
$\boxdot$ is proposed \emph{not} as $(E\Msum F)^{**}$ but directly as
$(E^*\fibsum F^*)^*$, mirroring the identity
$(f\,\square\,g)^*=f^*+g^*$ — which holds unconditionally at the function
level, unlike its converse — so that $\boxdot$ is automatically valued in
$\bar X$ without a separate nucleus-compatibility argument. The candidate
identity $E\boxdot F=(E\Msum F)^{**}$ is not claimed here; see
Milestone II, Question~\ref{q:representation}.

\subsection{Four units}

\begin{claim}[Units]\label{cl:units}
\[
1_\wedge:=X\times\Rp,\qquad
1_\vee:=X\times\{+\infty\},\qquad
1_\oplus:=X\times[0,+\infty],\qquad
1_\boxdot:=(\{0\}\times[0,+\infty])\cup\big((X\setminus\{0\})\times\{+\infty\}\big).
\]
\end{claim}

\emph{Motivation.} $1_\wedge,1_\vee$ mirror the constant functions
$-\infty$ (excluded from $\Rp$, hence $1_\wedge$ has no function-level
mirror under the properness convention) and $+\infty$. $1_\oplus$ mirrors
the zero function, the unit of pointwise addition. $1_\boxdot$ mirrors the
indicator $\delta_{\{0\}}$, the unit of infimal convolution, occupying
$X\times\{+\infty\}$ away from the origin — a shape forced by every
regepi's fiber being allowed to reach $+\infty$. \emph{Note:} these four
objects and their pairwise duality ($1_\wedge^*=1_\vee$,
$1_\oplus^*=1_\boxdot$) are claimed independently of which MALL connective
each is assigned to; if the $\otimes/\parr$ assignment in
Section~\ref{sec:components} is later reversed (Question~\ref{q:swap}),
the units themselves are unaffected — only their logical labels swap.

\section{Roadmap}\label{sec:milestones}

The three milestones are separated by \emph{what is being claimed about},
not by difficulty: Milestone I never mentions convexity or logic; Milestone
II interprets the algebra of Milestone I as convex analysis; Milestone III
interprets it as MALL. This separation lets Milestones I--II stand or fall
independently of whether the logic interpretation in Milestone III
survives.

\subsection*{Milestone I --- Intrinsic Definitions and Properties on Preepi}

Establishes the algebra of Claims~\ref{cl:closures}--\ref{cl:units} as a
self-contained mathematical structure on $\tilde X$, stated and proved
without reference to convex functions or to linear logic: extensivity
($E\subseteq E^{**}$), antitonicity, idempotence of $**$, regularizing
property ($E^*\in\bar X$ for every $E$), closure invariance
($E^*=U(E)^*=F(E)^*=C(E)^*=T(E)^*$), the bipolar theorem on $\bar X$ and its
general formula on arbitrary preepi; $(\bar X,\subseteq)$ a complete
lattice; $\oplus,\boxdot$ commutative monoids with the stated units,
monotone and distributive over arbitrary joins/meets as appropriate;
$(\cdot)^*$ an order-reversing involution on $\bar X$; De Morgan laws for
all four operators; duality of the four units. This is proposed as a
$*$-autonomous poset / Girard quantale purely as a poset-algebraic fact,
with no claim yet made about what it represents.

\subsection*{Milestone II --- Connection to Convex Optimization}

Interprets Milestone I as the theory of proper, convex, lower-semicontinuous
functions: $\bar X\leftrightarrow\Gamma_0(X)$ via $\epi$; the bipolar
theorem as Fenchel--Moreau, $f=f^{**}\iff f\in\Gamma_0(X)$; $\wedge,\vee$ as
$\max$ and closed $\min$; $\oplus,\boxdot$ as closed pointwise sum and
closed infimal convolution, in whichever direction Question~\ref{q:swap}
eventually resolves; the four units as $-\infty,+\infty,0,\delta_{\{0\}}$.

\begin{question}[Representation of $\boxdot$]\label{q:representation}
Does $E\boxdot F=(E\Msum F)^{**}$ on $\bar X$, i.e.\ does the polar-first
definition of Claim~\ref{cl:induced} agree with the naive
double-polar-of-Minkowski-sum definition used in earlier drafts? At the
function level this is the classical qualification question for exactness
of infimal convolution, e.g.\ under $\ri(\dom f)\cap\ri(\dom g)\ne\varnothing$.
\end{question}

\subsection*{Milestone III --- Connection to Logic}

This milestone claims nothing about preepi or convex functions; it records
the questions needed to justify reading Milestone I as a model of MALL.

\begin{question}[The $\otimes/\parr$ assignment]\label{q:swap}
Which of $\oplus,\boxdot$ should be identified with $\otimes$, and which
with $\parr$? Section~\ref{sec:components} adopts
$\otimes\leftrightarrow\oplus$, $\parr\leftrightarrow\boxdot$ on a
resource-intuition basis, but Girard's phase semantics defines $\otimes$
from a single ambient monoid, which on $\tilde X$ is $\Msum$ — favoring
the reverse assignment unless a genuinely different ambient structure
(e.g.\ a quantaloid indexed by $x\in X$, in the sense already attempted
and found non-associative in earlier drafts of this program) is used
instead.
\end{question}

\begin{question}[Phase semantics, made explicit]
Does the pairing $\beta$ and the polar of Claim~\ref{cl:polar} instantiate,
in the technical sense of Girard, a monoid acting on $\tilde X$ together
with a pole $\bot$, so that $E^*=\{a:\forall e\in E,\ e\cdot a\in\bot\}$
literally, and $\bar X=\Fact_\bot(\tilde X)$ in the standard phase-semantic
sense — rather than merely an algebra that happens to resemble one?
\end{question}

\begin{question}[Soundness]
Once Question~\ref{q:swap} is resolved, does the resulting correspondence
table validate each MALL inference rule: if the premises' semantic values
lie in the claimed facts, does the conclusion's value as well? Is the
semantic assignment $\llbracket\cdot\rrbracket$ well defined by structural
induction on formulas?
\end{question}

\begin{question}[Scope]
The program restricts itself to the poset/phase-semantic level of MALL:
no exponentials ($!,?$), no category of proofs, no cut elimination. Is this
restriction adequate for the intended application to convex analysis, or
does a categorical extension become necessary at some point?
\end{question}

\end{document}
